\section{Evaluation}
\label{sec:evaluation}

This section uses only the public benchmark families in the repository. Internal benchmark
families are intentionally excluded. The goal is not
to claim that the revised flow is always faster than the stock baseline. The checked-in
evidence shows the opposite in many cases because the manifest, injection, and rebuild stages
add cost. The evaluation instead answers two narrower questions:

\begin{enumerate}[leftmargin=1.5em]
  \item Do the public benchmark artifacts contain cases where the revised flow is judged more
  faithful than the stock asynchronous baseline?
  \item What cost is paid for that increase in accuracy?
\end{enumerate}

\subsection{Public Benchmark Scope}

The checked-in public report \texttt{check-src/benchmark.md} covers the five suites summarized
in \Cref{tab:public-suites}. Together they contribute 43 measured cases.

\begin{table}[t]
  \centering
  \small
  \begin{tabularx}{\linewidth}{@{}l c c X@{}}
    \toprule
    Suite & Cases & Headline-ready & Notes \\
    \midrule
    \texttt{icbmc-large} & 12 & yes & Larger i-CBMC-derived public cases. \\
    \texttt{icbmc-timeout-30m} & 4 & no & Public diagnostic cases kept outside the automatic
    suite. \\
    \texttt{icbmc-timeout-30m-blink-highmem} & 2 & no & Public high-memory blink cases. \\
    \texttt{icbmc-timeout-30m-logger-stable} & 4 & no & Public logger-focused diagnostic
    cases. \\
    \texttt{intabs-large} & 21 & yes & Larger public IntAbs-derived cases. \\
    \bottomrule
  \end{tabularx}
  \caption{Public benchmark suites used in this section.}
  \label{tab:public-suites}
\end{table}

\subsection{Public Correctness Summary}

The public subset has 43 measured cases. The correctness counts in
\Cref{tab:benchmark-summary} are recomputed for that public subset only.

\begin{table}[t]
  \centering
  \small
  \begin{tabularx}{\linewidth}{@{}l c X@{}}
    \toprule
    Benchmark verdict & Cases & Interpretation \\
    \midrule
    \texttt{same\_outcome\_comparable} & 17 & Stock and improved reached the same
    verification class. \\
    \texttt{improved\_correct} & 10 & The improved pipeline is audited as the more faithful
    model for the case. \\
    \texttt{no\_injection\_false} & 6 & The improved path reported no injection candidates
    even though state written by an interleaving function can still affect the main-reachable
    state. \\
    \texttt{needs\_manual\_review} & 3 & Output mismatch exists, but the current audit is not
    strong enough to choose a correct variant. \\
    \texttt{pending\_manual\_review} & 7 & Mismatch exists and no verdict rule has yet
    classified the case. \\
    \bottomrule
  \end{tabularx}
  \caption{Correctness summary for the public benchmark subset only.}
  \label{tab:benchmark-summary}
\end{table}

Two points stand out. First, the improved pipeline is not merely different from the stock
baseline; it is often audited as more faithful. Second, the benchmark report explicitly keeps
manual-review categories instead of silently forcing every mismatch into a win/loss label. That
is important engineering hygiene for an evolving verification branch.

\subsection{Representative Public Cases}

Three public cases illustrate the different ways in which the revised flow changes the
verification result.

\paragraph{\texttt{icbmc-timeout-30m-logger-stable / logger2-bug-conc-cprover}.}
The stock baseline reports \emph{VERIFICATION FAILED} in 56\,233 ms with 298.2 MB peak RSS,
while the improved pipeline reports \emph{VERIFICATION SUCCESSFUL} in 2\,017 ms with
16.6 MB peak RSS. The checked-in audit explains the difference directly: the stock async
baseline exposes a non-atomic asynchronous interleaving in the logger update, whereas the
improved pipeline checks the extracted interruption entry as one atomic region. This is a
clear public example of unnecessary interleavings producing a spurious bug.

\paragraph{\texttt{intabs-large / i8xx-tco-1}.}
The stock baseline reports \emph{VERIFICATION FAILED} in 1202 ms with 13.1 MB peak RSS, while
the improved pipeline reports \emph{VERIFICATION SUCCESSFUL} in 3395 ms with 15.3 MB peak RSS.
The benchmark audit again points to unnecessary interleavings in the stock model. The stock
run allows a non-atomic asynchronous interleaving, whereas the improved pipeline keeps the
interruption entry atomic and narrows the candidate set to three guarded interleavings. This
case shows the main accuracy gain of the revised flow even though the full measured time grows
from 1202 ms to 3395 ms.

\paragraph{\texttt{intabs-large / sc520wdt-2}.}
The stock baseline reports \emph{VERIFICATION SUCCESSFUL} in 1172 ms with 41.0 MB peak RSS.
The improved pipeline reports \emph{VERIFICATION FAILED} in 3559 ms with 16.7 MB peak RSS.
Here the improved pipeline injects closer/writer interleavings and reaches a violated
property that the stock baseline proves away. This is the opposite public pattern from the
previous two cases: the revised flow is slower, but it exposes an error trace that is not
reached by the stock model. The full measured time rises from 1172 ms to 3559 ms, but the
peak RSS drops from 41.0 MB to 16.7 MB.

Taken together, these public cases show the main cost of the revised flow. Even when the
verification phase itself is not dramatically slower, the extra manifest, injection, and
rebuild phases raise the full measured time. The public evidence therefore supports a clear
trade-off: the revised flow aims at a more accurate verification result, not at a faster one.

For the use case of this report, that trade is acceptable. The objective is not to outperform
the stock baseline on total wall-clock time, but to reduce unnecessary interleavings and make
bounded model checking more accurate for OSEK/VDX applications.
